\documentclass[10pt, a4paper]{article}

\usepackage[a4paper, top=2cm, bottom=2cm, left=2.2cm, right=2.2cm]{geometry}
\usepackage{amsmath, amssymb}
\usepackage{booktabs}
\usepackage{array}
\usepackage{tabularx}
\usepackage{multirow}
\usepackage{hyperref}
\usepackage{enumitem}
\usepackage{microtype}
\usepackage{parskip}
\usepackage{titlesec}
\usepackage{caption}
\usepackage{float}

\setlength{\parskip}{4pt}
\setlength{\parindent}{0pt}

\hypersetup{colorlinks=true, linkcolor=black, citecolor=black, urlcolor=black}

\titleformat{\section}{\normalsize\bfseries}{}{0em}{}[\vspace{-4pt}\rule{\linewidth}{0.8pt}]
\titleformat{\subsection}{\normalsize\bfseries}{\thesubsection}{0.8em}{}
\titleformat{\subsubsection}{\normalsize\itshape}{\thesubsubsection}{0.8em}{}
\titlespacing{\section}{0pt}{10pt}{4pt}
\titlespacing{\subsection}{0pt}{6pt}{2pt}
\titlespacing{\subsubsection}{0pt}{4pt}{2pt}

\captionsetup{font=small, skip=4pt}

\begin{document}

% ── Header ───────────────────────────────────────────────────────────────────
\begin{center}
  {\large\textbf{MIS: Difficulty Scoring as a Pre-Augmentation Diagnostic
    for Agricultural Object Detection}}\\[4pt]
  Huy Khanh Tran \quad
  Graduate School of Integrated Science and Technology,
  Shizuoka University\\[2pt]
  {\small Research Status Report --- May 2026 --- Target: GCCE 2026 (2 pages)}
\end{center}
\rule{\linewidth}{1pt}
\vspace{2pt}

% ── Summary ──────────────────────────────────────────────────────────────────
\textbf{One-sentence claim.}
MIS provides structural diagnostics --- computable from training-set inference
alone --- that are consistent with observed copy-paste augmentation failure
and that identify the failure conditions before augmentation is conducted.

\textbf{Why the claim changed from the original draft.}
The original aim was to show that MIS-guided copy-paste improves AP.
Three datasets (MinneApple, GWHD, VisDrone) showed uniform failure.
The original positive result on MinneApple ($\Delta\text{AP}=+0.013$
for Score\,CP vs.\ Random\,CP) was an artifact of a contaminated validation
split and does not hold on clean data.
Rather than search for a positive result, the research is reframed:
the consistent failure is itself informative, and MIS correctly identifies
the structural reasons for that failure.

% ══════════════════════════════════════════════════════════════════════════════
\section{Methodology}
% ══════════════════════════════════════════════════════════════════════════════

\subsection{MIS Scoring}

Predictions are generated at low confidence threshold $\tau_{\text{conf}}=0.001$
(to capture near-miss detections) and matched to ground truth at
$\tau_{\text{IoU}}=0.5$.
Let $\mathcal{P}(g)$ be matched predictions for object $g$.
The \textbf{object difficulty score} is:
\begin{equation}
  S_{\text{obj}}(g) = \begin{cases}
    \min_{p \in \mathcal{P}(g)}
      \bigl[\alpha(1 - c_p) + \beta(1 - u_p)\bigr],
      & \mathcal{P}(g) \neq \varnothing \\
    \alpha + \beta = 1.0, & \text{missed object}
  \end{cases}
  \label{eq:sobj}
\end{equation}
with $\alpha=\beta=0.5$, $c_p=\text{conf}(p)$, $u_p=\text{IoU}(p,g)$.
Taking the minimum is necessary because the low $\tau_{\text{conf}}$
generates many candidate boxes; minimum selects the best-matching one.

The \textbf{image difficulty score} is:
\begin{equation}
  S_{\text{img}}(I) = w_1 \cdot \frac{1}{|G(I)|}
    \sum_{g \in G(I)} S_{\text{obj}}(g)
    + w_2 \cdot \text{FPrate}(I)
  \label{eq:simg}
\end{equation}
$w_1 = 1$ (fixed). $w_2$ is calibrated automatically by grid search to
minimize $|r(S_{\text{img}},\text{miss\_rate})-r(S_{\text{img}},\text{FP\_rate})|$,
balancing both error signals.
FP rate is computed at the optimal operating threshold found from validation.

\subsection{Dataset Splits}

\textbf{MinneApple.}
The dataset contains images from only 10 unique apple trees.
A random image-level split causes images from the same tree to appear in
both training and validation, inflating validation AP through tree-specific
memorization (initial contaminated baseline: AP$=0.439$, discarded).
We use a \textbf{tree-identity-clean split}: validation uses 30 images
from trees not in training ($\approx$3 images per held-out tree).
Test set is the official MinneApple test split.
Training: 200 epochs max, early stopping patience 50 epochs.

\textbf{MinneApple size thresholds.}
COCO thresholds (small $\leq32$\,px, medium $\leq96$\,px) do not suit
MinneApple at $720{\times}1280$ resolution.
We calibrate to the actual training distribution:

\begin{center}
\small
\begin{tabular}{lll}
\toprule
\textbf{Category} & \textbf{Criterion} & \textbf{Train \%} \\
\midrule
Very Small & area $\leq400$\,px$^2$ ($<20{\times}20$\,px) & 23.08\% \\
Small      & $400<$ area $\leq1024$\,px$^2$ (20--32\,px)  & 50.53\% \\
Medium     & area $>1024$\,px$^2$ ($>32{\times}32$\,px)   & 26.39\% \\
\bottomrule
\end{tabular}
\end{center}

\subsection{Augmentation Setup}

A YOLOv8s baseline is trained from COCO pretrained weights.
Inference on the full training set yields scores via Eq.\,\eqref{eq:sobj}--\eqref{eq:simg}.
Training set extended by 30\% (189 synthetic images added to 536).
Each synthetic image: 2--8 pasted objects, scale $[0.9,1.1]$,
rotation $\pm15^\circ$, no blending, \textbf{same-image constraint}
(objects pasted into the image they come from, to preserve photometric
and contextual consistency, following Kisantal et al.\ \cite{kisantal2019}).

All images and objects are ranked by score and partitioned into three
equal tertiles. Five conditions compared (3 seeds each):

\begin{center}
\small
\begin{tabular}{lll}
\toprule
\textbf{Condition} & \textbf{Image pool} & \textbf{Object pool} \\
\midrule
Random        & Uniform random     & Uniform random \\
Score High    & Top tertile        & Top tertile \\
Score Medium  & Middle tertile     & Middle tertile \\
Score Low     & Bottom tertile     & Bottom tertile \\
\bottomrule
\end{tabular}
\end{center}

% ══════════════════════════════════════════════════════════════════════════════
\section{Score Validation}
% ══════════════════════════════════════════════════════════════════════════════

\subsection{Image-Level Correlation}

\begin{table}[H]
\centering
\small
\caption{Pearson correlation of $S_{\text{img}}$ with per-image error rates.}
\begin{tabular}{lcc}
\toprule
\textbf{Dataset} & $r(\text{score, miss rate})$ & $r(\text{score, FP rate})$ \\
\midrule
MinneApple & 0.698 & 0.688 \\
GWHD       & 0.816 & 0.818 \\
\bottomrule
\end{tabular}
\end{table}

Moderate-to-strong correlations on both datasets.
GWHD is higher because its difficulty is structured at the domain level
(cleaner signal), while MinneApple has heterogeneous within-tree variation.

\subsection{Object-Size Validity: MinneApple}

\begin{table}[H]
\centering
\small
\caption{Mean $S_{\text{obj}}$ and baseline AP by size category (MinneApple,
  consistent across all 3 seeds). Spearman $r_s=-0.82$ between size rank
  and difficulty rank confirms strong monotonic inverse relationship.}
\begin{tabular}{lcccc}
\toprule
\textbf{Category} & \textbf{AP} & \textbf{Mean $S_{\text{obj}}$}
  & \textbf{Std} & \textbf{In top-score tertile} \\
\midrule
Very Small ($<$20\,px) & 0.071 & 0.1524 & 0.1433
  & \textbf{64.39\% (5138 objects)} \\
Small (20--32\,px)     & 0.270 & 0.0755 & 0.0475
  & 33.88\% (2704 objects) \\
Medium ($>$32\,px)     & 0.493 & 0.0526 & 0.0184
  & \textbf{1.73\% (138 objects)} \\
\bottomrule
\end{tabular}
\end{table}

\textbf{Key finding: 98.27\% of high-difficulty objects are in size
categories where baseline AP $\leq$ 0.27.}
Only 1.73\% of the objects MIS identifies as hardest fall in the medium
category where the model has genuine learning headroom (AP$=0.49$).
This means score-guided augmentation concentrating on hard objects
will predominantly paste near-undetectable instances regardless of
which difficulty tier is selected.

\subsection{Domain-Level Validity: GWHD}

\begin{table}[H]
\centering
\small
\caption{GWHD domain analysis: MIS median difficulty vs.\ AP across
  18 domains. Pearson $r=-0.905$, Spearman $r_s=-0.899$.
  MIS ranks domain difficulty without domain labels.}
\begin{tabular}{lcc||lcc}
\toprule
\textbf{Domain} & \textbf{Median $S$} & \textbf{AP}
  & \textbf{Domain} & \textbf{Median $S$} & \textbf{AP} \\
\midrule
Arvalis\_10     & 0.200 & 0.744 & Arvalis\_8      & 0.342 & 0.546 \\
Rres\_1         & 0.227 & 0.658 & Arvalis\_5      & 0.342 & 0.542 \\
ETHZ\_1         & 0.268 & 0.643 & Arvalis\_9      & 0.372 & 0.520 \\
NMBU\_2         & 0.299 & 0.639 & Arvalis\_12     & 0.395 & 0.472 \\
Inrae\_1        & 0.264 & 0.633 & Arvalis\_3      & 0.350 & 0.470 \\
Arvalis\_11     & 0.358 & 0.606 & Arvalis\_7      & 0.407 & 0.463 \\
NMBU\_1         & 0.335 & 0.601 & Arvalis\_2      & 0.386 & 0.454 \\
Arvalis\_6      & 0.289 & 0.600 & Arvalis\_1      & 0.400 & 0.452 \\
ULiège-GxABT\_1 & 0.380 & 0.597 & Arvalis\_4      & 0.475 & 0.387 \\
\bottomrule
\end{tabular}
\end{table}

% ══════════════════════════════════════════════════════════════════════════════
\section{Augmentation Study}
% ══════════════════════════════════════════════════════════════════════════════

\subsection{Object Size Distribution Across Conditions}

\begin{table}[H]
\centering
\small
\caption{Object scale distribution per augmentation condition.
  \emph{Gap} = medium-object \% minus test-set target (44.70\%).
  All gaps are negative: no condition bridges the train--test size mismatch.}
\begin{tabular}{lrrrrc}
\toprule
\textbf{Condition} & \textbf{V.Small} & \textbf{Small}
  & \textbf{Medium} & \textbf{Gap} & \textbf{Obj/Img} \\
\midrule
Test target       & 17.17\% & 38.14\% & \textbf{44.70\%}
  & 0\,pp & 37.11 \\
\midrule
Baseline (train)  & 23.08\% & 50.53\% & 26.39\% & $-$18.31\,pp & 42.57 \\
Random CP         & 23.52\% & 49.98\% & 26.50\% & $-$18.20\,pp & 44.03 \\
Score High        & 26.48\% & 49.59\% & 23.92\% & $-$20.78\,pp & 43.85 \\
Score Medium      & 22.56\% & 50.80\% & 26.64\% & $-$18.06\,pp & 43.23 \\
\textbf{Score Low}& \textbf{19.99\%} & 48.56\% & \textbf{31.45\%}
  & $-$13.25\,pp & 42.38 \\
\bottomrule
\multicolumn{6}{l}{\small Score High \emph{widens} the gap.
  Score Low achieves the largest shift toward test, yet still fails (see Table 4).}
\end{tabular}
\end{table}

The maximum density increase across all conditions is \textbf{3.4\%}
(42.57 $\to$ 44.03 objects per image).
Kisantal et al.\ achieved density increases of 30--60\% starting from
3--5 objects per image.
The mechanism that produced gains in their setting does not operate
at MinneApple's baseline density.

\subsection{Detection Results}

\begin{table}[H]
\centering
\small
\caption{3-seed mean AP with standard deviation.
  All $\Delta$AP values are not statistically significant (ns):
  all differences $<1\sigma$ of baseline.}
\begin{tabular}{lcrrrr}
\toprule
\textbf{Condition} & \textbf{AP $\pm\sigma$} & $\Delta$\textbf{AP}
  & \textbf{AP$_{50}$} & \textbf{AP\,vs} & \textbf{AP\,s} \\
\midrule
Baseline      & \textbf{0.353 $\pm$ 0.005} & ---
  & 0.648 & 0.071 & 0.270 \\
Random CP     & 0.346 $\pm$ 0.005 & $-$0.007 (ns) & 0.637 & 0.069 & 0.266 \\
Score High    & 0.350 $\pm$ 0.007 & $-$0.003 (ns) & 0.645 & 0.069 & 0.271 \\
Score Medium  & 0.348 $\pm$ 0.006 & $-$0.005 (ns) & 0.643 & 0.072 & 0.272 \\
Score Low     & 0.348 $\pm$ 0.003 & $-$0.005 (ns) & 0.643 & 0.068 & 0.271 \\
\bottomrule
\multicolumn{6}{l}{\small AP\,vs = AP very small ($<$20\,px).
  AP\,s = AP small (20--32\,px).}
\end{tabular}
\end{table}

\textbf{Score High} selects 64.39\% very small objects (AP$=0.07$),
widens the distribution gap to 20.78\,pp, and produces no improvement.
Confirms that hard objects in MinneApple are near-undetectable
regardless of training exposure.

\textbf{Score Low} (most important result) selects predominantly
medium-sized objects, achieves the largest distribution shift
(gap narrows to 13.25\,pp), yet still produces no improvement.
This eliminates selection strategy as the explanation for failure:
even the condition best positioned to succeed cannot overcome
a 3.4\% density increase on a dataset already averaging 42.57 objects
per image.

\subsection{Feature Space Analysis}

Feature representations extracted from YOLOv8s P5 layer,
projected via PCA (PC1: 24.4\%, PC2: 15.4\% variance).

\begin{center}
\small
\begin{tabular}{lrr}
\toprule
\textbf{Split} & \textbf{PC1 centroid} & \textbf{PC2 centroid} \\
\midrule
Original train  & $+4.83$ & $+1.30$ \\
Augmented       & $+3.72$ & $-3.83$ \\
Test            & $-11.24$ & $-0.35$ \\
\bottomrule
\end{tabular}
\end{center}

Augmented images remain co-located with training data
($\Delta\text{PC1}\approx 1.1$) while both are separated from the
test distribution by $\Delta\text{PC1}\approx 16$.
\textbf{Copy-paste does not move the training feature distribution
toward the test distribution.}

% ══════════════════════════════════════════════════════════════════════════════
\section{Diagnostic Framework}
% ══════════════════════════════════════════════════════════════════════════════

The purpose of this section is to answer the question:
\textbf{given MIS scores on a dataset, how do you decide whether to
apply copy-paste augmentation?}

The answer is a three-criterion checklist.
Each criterion is checkable from training-set inference alone, before
any augmentation is run.

\subsection{The Three Criteria}

\textbf{Criterion 1: AP of hard objects (resolution ceiling).}\\
Compute $S_{\text{obj}}$ for all training objects.
Find the top-difficulty tertile (hardest third by score).
Identify which size/domain category those objects predominantly belong to.
Look up the baseline AP for that category.

\begin{center}
\small
\begin{tabular}{lll}
\toprule
\textbf{Condition} & \textbf{Signal} & \textbf{Action} \\
\midrule
AP of hard-object category $<$ 0.15
  & Hard objects near-undetectable.
  & \textbf{Do not use copy-paste.} \\
  & Pasting more copies adds no signal. & \\
AP of hard-object category $\geq$ 0.15
  & Hard objects partially learnable.
  & Proceed to Criterion 2. \\
\bottomrule
\end{tabular}
\end{center}

\textit{MinneApple result:} AP$_\text{very small}=0.071$, and 64.39\% of
hard objects are very small. \textbf{Criterion 1 signals: do not use copy-paste.}

\medskip
\textbf{Criterion 2: Domain concentration (covariate shift).}\\
For each image in the top-difficulty tertile, record which domain it
comes from. If the majority come from domain-minority sources:

\begin{center}
\small
\begin{tabular}{ll}
\toprule
\textbf{Condition} & \textbf{Action} \\
\midrule
Hard images concentrated in minority domains
  & Score-guided selection amplifies imbalance. \\
  & \textbf{Do not use score-guided copy-paste.} \\
Hard images spread across domains
  & Proceed to Criterion 3. \\
\bottomrule
\end{tabular}
\end{center}

\textit{GWHD result:} Hard images concentrate in the Arvalis minority
domains (lowest AP). \textbf{Criterion 2 signals: do not use score-guided
copy-paste.}

\medskip
\textbf{Criterion 3: Augmentation density headroom.}\\
Compute: proposed density increase $= \frac{\text{pasted objects/image} \times \text{augmentation ratio}}{\text{mean objects/image}}$.

\begin{center}
\small
\begin{tabular}{ll}
\toprule
\textbf{Condition} & \textbf{Action} \\
\midrule
Density increase $<$ 10\%
  & Augmentation too small to shift distribution. \\
  & \textbf{Increase augmentation ratio or reduce baseline density.} \\
Density increase $\geq$ 10\%
  & Distribution shift is viable. Copy-paste may help. \\
\bottomrule
\end{tabular}
\end{center}

\textit{MinneApple result:} $(5\ \text{obj} \times 0.35) / 42.57 = 4.1\%$
density increase. \textbf{Criterion 3 signals: insufficient headroom.}

\subsection{Summary Decision Table}

\begin{table}[H]
\centering
\small
\caption{Applying the three-criterion diagnostic to both datasets.
  All three criteria signal failure on both.
  All confirmed by experiment (Table 4).}
\begin{tabular}{lllll}
\toprule
\textbf{Criterion} & \textbf{MinneApple} & \textbf{GWHD}
  & \textbf{Prediction} & \textbf{Observed} \\
\midrule
1. AP of hard objects
  & 0.071 $\to$ \textbf{Fail}
  & Mixed (domain-dependent)
  & Fail & Fail\,$\checkmark$ \\
2. Domain concentration
  & N/A (single domain)
  & Minority-concentrated $\to$ \textbf{Fail}
  & Fail & Fail\,$\checkmark$ \\
3. Density headroom
  & 3.4\% $\to$ \textbf{Fail}
  & Similar $\to$ \textbf{Fail}
  & Fail & Fail\,$\checkmark$ \\
\midrule
\multicolumn{2}{l}{\textbf{Overall}} & & \textbf{Fail} & \textbf{Fail\,$\checkmark$} \\
\bottomrule
\end{tabular}
\end{table}

\subsection{When Copy-Paste Is Likely to Help}

The framework also implies the positive conditions --- the dataset
properties where copy-paste is worth attempting:

\begin{itemize}[nosep, leftmargin=*]
\item AP of hard objects $\geq 0.15$ (learnable, not yet saturated),
\item High-difficulty images spread across domains (no concentration),
\item Density increase $\geq 10\%$ (sparse dataset, augmentation moves
  the distribution).
\end{itemize}

These conditions are consistent with Kisantal et al.'s success on COCO
(sparse images, 30--60\% density increase, large learnable objects).
On MinneApple and GWHD, none of these conditions hold, which explains
the uniform failure.

\subsection{Limitations of the Framework}

The three criteria were derived from the same datasets used to validate
them, so this is a retrospective characterization, not a prospective
prediction tested on an unseen dataset.
The thresholds (AP $<0.15$, density $<10\%$) are operationally
reasonable but not formally derived.
A prospective test on a new dataset would strengthen the framework
substantially.

% ══════════════════════════════════════════════════════════════════════════════
\section{Proposed Paper Structure (GCCE, 2 pages)}
% ══════════════════════════════════════════════════════════════════════════════

\begin{center}
\small
\begin{tabular}{llp{10cm}}
\toprule
\textbf{Section} & \textbf{$\approx$Words} & \textbf{Content and main point} \\
\midrule
Abstract      & 100 & MIS scorer. 5 conditions. 64.39\% hard = very small.
  18.31\,pp mismatch. Decision rule. No positive AP claimed. \\
I. Intro      & 160 & Copy-paste fails on agricultural data.
  No diagnostic exists.
  Two contributions: (1) MIS scorer, (2) structural diagnostic study. \\
II. Related   & 100 & OHEM/Focal Loss (training-time, different purpose).
  Kisantal (anchor-based, mechanism does not transfer to YOLOv8s).
  MIS: post-hoc, architecture-agnostic. \\
III. Method   & 200 & Eq.~1--2 with correct $\tau$ values.
  Five conditions. Same-image justification.
  Custom thresholds with justification. \\
IV. Diagnostic& 200 & Table: Pearson. Table: size rankings + 64.39\% finding.
  GWHD $r=-0.905$. 18\,pp mismatch.
  Three-criterion decision table. \\
V. Augmentation & 250 & Combined AP + distribution table.
  Score Low analysis (most informative result).
  PCA figure. All deltas marked (ns). \\
VI. Conclusion & 100 & Structural failure confirmed.
  Decision rule in 3 steps.
  Future: generative augmentation targeting the learnable zone. \\
\bottomrule
\end{tabular}
\end{center}

% ══════════════════════════════════════════════════════════════════════════════
\section{Known Errors to Fix Before Submission}
% ══════════════════════════════════════════════════════════════════════════════

\begin{enumerate}[nosep, leftmargin=*]
\item \textbf{Verify GWHD split.}
  Confirm GWHD used a domain-based split (not random image-level).
  If contaminated, retrain or limit claims to MinneApple only.
\item \textbf{Rename ``boundary sampling.''} 
  The implementation is tertile filtering, not class-boundary sampling
  in the ML sense. Use ``difficulty-tier selection'' throughout.
\item \textbf{Do not write ``predicts.''}
  The framework is retrospectively validated on the same datasets
  used to derive it.
  Write ``identifies structural conditions consistent with failure''
  rather than ``predicts failure.''
\item \textbf{Change ``strongly positive''} to ``moderate-to-strong
  ($r=0.698$--$0.818$)'' for Pearson correlations.
  $r=0.698$ is moderate, not strong.
\item \textbf{Fix all typos} from original draft before submission.
\end{enumerate}

% ── References ───────────────────────────────────────────────────────────────
\begin{thebibliography}{9}

\bibitem{kisantal2019}
M.~Kisantal, Z.~Wojna, J.~Murawski, J.~Naruniec, and K.~Cho,
``Augmentation for small object detection,''
\textit{arXiv:1902.07296}, 2019.

\bibitem{ghiasi2021}
G.~Ghiasi et al.,
``Simple copy-paste is a strong data augmentation method for instance
segmentation,''
in \textit{Proc.\ IEEE/CVF CVPR}, 2021, pp.\ 2918--2928.

\bibitem{minneapple}
N.~Hani, P.~Roy, and B.~Bhanu,
``MinneApple: A benchmark dataset for apple detection and segmentation,''
\textit{IEEE RA-L}, vol.~5, no.~2, pp.\ 852--858, 2020.

\bibitem{gwhd}
E.~David et al.,
``Global Wheat Head Detection (GWHD) dataset,''
\textit{Plant Phenomics}, 2021.

\bibitem{ohem}
A.~Shrivastava, A.~Gupta, and R.~Girshick,
``Training region-based object detectors with online hard example mining,''
in \textit{Proc.\ IEEE CVPR}, 2016.

\bibitem{focal}
T.-Y.~Lin et al.,
``Focal loss for dense object detection,''
in \textit{Proc.\ IEEE ICCV}, 2017.

\bibitem{yolov8}
G.~Jocher, A.~Chaurasia, and J.~Qiu,
``YOLO by Ultralytics,'' v8.0.0, 2023.

\end{thebibliography}

\end{document}